\section{Resultados}

Esta seção apresenta os resultados obtidos no estudo de caso \textit{poster}. A descrição do treinamento, da preparação das variantes e da integração em Unity é mantida na seção de Implementação; aqui são apresentados os dados derivados dos registros experimentais, sua interpretação delimitada e as limitações correspondentes. As tabelas foram consolidadas a partir de arquivos de evidência e de cinco execuções independentes por variante quando essa série está indicada.

Para tornar a força de cada afirmação explícita, valores numéricos e configurações são tratados como fatos medidos; descrições de imagens e anotações realizadas durante o uso do headset são tratados como observações qualitativas; e as relações entre custo, fidelidade e adequação para VR são apresentadas como interpretações restritas ao protocolo. Nenhuma dessas categorias é usada para extrapolar o resultado para cenas, poses ou dispositivos não avaliados.

\subsection{Escopo da avaliação}

As métricas offline foram calculadas sobre as mesmas 13 imagens de teste retidas do conjunto \textit{poster}. A execução \textit{on-device} ocorreu no Meta Quest 3S, em condição de referência estática, com pose completa fixa após o rastreamento. Para evitar comparar observações com forças de evidência diferentes, as variantes \textit{baseline}, \textit{top-k} 100k e \textit{top-k} 50k formam a série comparativa de cinco execuções. A variante \textit{splatfacto-big} é apresentada separadamente como diagnóstico estático de uma execução.

Para consolidar os resultados, as medições originais foram mantidas separadas das tabelas e dos gráficos, e foi conferida a correspondência entre as representações avaliadas e as configurações de cada execução.

\subsection{Qualidade offline e complexidade das representações}

O Quadro \ref{qua:resultados_offline} apresenta as métricas das representações no mesmo conjunto de imagens retidas. Como fatos medidos, a regularização de escala teve menos gaussianas que a \textit{baseline}, menor PSNR e SSIM e maior LPIPS. A configuração \textit{splatfacto-big} teve SSIM mais alto e LPIPS mais baixo, embora PSNR ligeiramente inferior ao da \textit{baseline}, e seu arquivo PLY foi maior.

\begin{quadro}[H]
\centering
\footnotesize
\caption{Qualidade offline e complexidade das representações no estudo de caso \textit{poster}.}
\label{qua:resultados_offline}
\renewcommand{\arraystretch}{1.2}
\begin{tabularx}{\textwidth}{|p{2.5cm}|C{1.75cm}|C{1.55cm}|C{1.55cm}|C{1.8cm}|C{1.7cm}|X|}
\hline
\textbf{Representação} & \textbf{PSNR} & \textbf{SSIM} & \textbf{LPIPS} & \textbf{Gaussianas} & \textbf{PLY (MB)} & \textbf{Tempo de treinamento} \\ \hline
\textit{baseline} & 32,989 & 0,95363 & 0,10495 & 195.760 & 48,55 & 1.357,7 s \\ \hline
\textit{top-k} 100k & 22,716 & 0,80851 & 0,31945 & 100.000 & 24,80 & -- \\ \hline
\textit{top-k} 50k & 17,514 & 0,60303 & 0,51890 & 50.000 & 12,40 & -- \\ \hline
Regularização de escala & 32,945 & 0,95310 & 0,11010 & 181.163 & 44,93 & 1.281,5 s \\ \hline
\textit{splatfacto-big} & 32,764 & 0,95407 & 0,08538 & 470.962 & 116,80 & 1.913,5 s \\ \hline
\end{tabularx}
\legend{Fonte: Elaborado pelo autor.}
\end{quadro}

As variantes por poda determinística \textit{top-k} foram derivadas da \textit{baseline}; portanto, não representam novo treinamento. Para avaliá-las, a configuração \textit{splatfacto} da \textit{baseline} foi carregada e seus parâmetros gaussianos foram substituídos pelos PLYs podados, preservando as mesmas 13 poses, imagens e resolução de teste. A reavaliação da \textit{baseline} reproduziu as médias registradas anteriormente. Como fatos medidos, a poda para 100 mil gaussianas resultou em PSNR 22,716, SSIM 0,80851 e LPIPS 0,31945; a poda para 50 mil gaussianas resultou em PSNR 17,514, SSIM 0,60303 e LPIPS 0,51890. A interpretação limitada é que, nesse conjunto de teste e nesse critério de seleção por opacidade, a redução do PLY não preservou as métricas offline da \textit{baseline}. A regularização de escala não foi levada ao Quest porque seus resultados offline não justificaram ampliar a série de execução no dispositivo.

\subsection{Desempenho em execução nativa no Quest}

O Quadro \ref{qua:resultados_quest} apresenta a série comparativa de cinco execuções por variante. O FPS principal é calculado a partir do tempo de quadro registrado pela aplicação; o FPS médio do OVR Metrics é apresentado como contador complementar, e não como medida intercambiável. Na série retida, as variantes com menos gaussianas apresentaram menores tempos médios de quadro e maiores FPS derivados.

\begin{quadro}[H]
\centering
\footnotesize
\caption{Desempenho no Meta Quest 3S em condição de referência estática, com cinco execuções por variante.}
\label{qua:resultados_quest}
\renewcommand{\arraystretch}{1.2}
\begin{tabularx}{\textwidth}{|p{2.6cm}|C{1.8cm}|C{1.8cm}|C{1.8cm}|C{1.9cm}|X|}
\hline
\textbf{Representação} & \textbf{Tempo de quadro (ms)} & \textbf{FPS da aplicação} & \textbf{FPS OVR} & \textbf{GPU app (ms)} & \textbf{Quadros obsoletos} \\ \hline
\textit{baseline} & 45,05 & 22,20 & 25,14 & 39,84 & 47,93 \\ \hline
\textit{top-k} 100k & 26,30 & 38,02 & 40,49 & 23,07 & 33,76 \\ \hline
\textit{top-k} 50k & 17,13 & 58,37 & 59,64 & 15,11 & 13,41 \\ \hline
\end{tabularx}
\legend{Fonte: Elaborado pelo autor.}
\end{quadro}

Como fatos medidos, a \textit{baseline} registrou 22,20 FPS, a poda 100k registrou 38,02 FPS e a poda 50k registrou 58,37 FPS, com o menor tempo médio de quadro da série. Nenhuma média atingiu a referência de 90 FPS definida no protocolo. Em interpretação restrita a essas três variantes e à pose fixa, a poda 50k oferece o menor custo de renderização da série, mas não constitui uma seleção definitiva para VR: seus valores de PSNR e SSIM foram os menores da série, enquanto o LPIPS foi o maior e as capturas qualitativas mostraram mais fragmentação. A condição estacionária não permite inferir fluidez, conforto ou estabilidade durante movimento de cabeça.

\begin{figure}[H]
\centering
\fbox{\includegraphics[width=0.92\textwidth]{figuras/resultados/quest_tradeoff_v02.png}}
\caption{Relação descritiva entre número de gaussianas e FPS derivado do tempo de quadro na execução nativa.}
\label{fig:quest_tradeoff}
\legend{Fonte: Elaborado pelo autor.}
\end{figure}

As capturas estéreo da Figura \ref{fig:quest_capturas_estereo} oferecem evidência qualitativa complementar sob a mesma condição de referência. Na inspeção das três imagens, a estrutura global da cadeira permaneceu reconhecível, enquanto a fragmentação e a perda de detalhes no piso e no fundo foram mais aparentes nas versões podadas. Essa é uma observação das imagens selecionadas, e não uma métrica de fidelidade. Por serem screenshots espaçados de uma pose fixa, elas não medem cintilação quadro a quadro, conforto ou qualidade óptica do visor.

\begin{figure}[H]
\centering
\fbox{\includegraphics[height=0.78\textheight,keepaspectratio]{figuras/resultados/quest_stereo_comparison_v01.png}}
\caption{Capturas estéreo da execução nativa no Quest, na mesma cena e pose fixa: \textit{baseline}, poda 100k e poda 50k, de cima para baixo.}
\label{fig:quest_capturas_estereo}
\legend{Fonte: Elaborado pelo autor.}
\end{figure}

\subsection{Diagnóstico de maior capacidade}

A variante \textit{splatfacto-big}, com 470.962 gaussianas e 116,80 MB, foi executada uma vez na mesma condição de pose fixa. O tempo de quadro medido pela aplicação foi 65,62 ms, correspondente a 15,24 FPS. Esse resultado não foi agregado à série de cinco execuções e não possui os contadores OVR disponíveis para comparação. No recorte dessa única execução, o FPS derivado ficou abaixo da referência de desempenho do protocolo; ele não permite estimar variabilidade entre execuções nem generalizar o efeito do aumento de complexidade.

Durante a observação qualitativa no headset, foram identificados nessa variante fundo com mais detalhes e menor distorção, cor aparentemente correta, menos cintilação e cadeira visualmente correta. Trata-se de observação restrita à cena \textit{poster}, à pose fixa e a uma execução; não constitui métrica perceptual nem permite concluir que a variante é superior em outras poses, cenas ou dispositivos.

\subsection{Contingência diagnóstica por PCVR}

Como contingência distinta da execução \textit{standalone}, foram realizadas execuções em Windows com transmissão para o Meta Quest 3S via Meta Horizon Link. O \textit{build} utilizou Direct3D12 e IL2CPP, com taxa de atualização reportada de 72 Hz. O Quadro \ref{qua:resultados_pcvr} apresenta uma única execução tecnicamente válida por variante; portanto, os valores são descritivos e não constituem repetição independente nem são comparados estatisticamente com a série Android.

Para essa avaliação, foi utilizado um computador com processador AMD Ryzen 7 5700G, oito núcleos e 16 processadores lógicos, GPU NVIDIA GeForce RTX 3060 e Windows 11 Pro de 64 bits. Foram mantidas a cena, as representações e a condição de pose fixa utilizadas na avaliação nativa. No PCVR, o computador renderizou a aplicação, enquanto o Meta Quest 3S recebeu as imagens transmitidas e forneceu o rastreamento. Essa configuração permitiu observar o comportamento das mesmas representações em uma plataforma de processamento distinta.

\begin{quadro}[H]
\centering
\footnotesize
\caption{Diagnóstico PCVR via Meta Horizon Link em pose fixa.}
\label{qua:resultados_pcvr}
\renewcommand{\arraystretch}{1.2}
\begin{tabularx}{\textwidth}{|p{3.0cm}|C{1.7cm}|C{1.9cm}|C{1.9cm}|C{1.8cm}|X|}
\hline
\textbf{Representação} & \textbf{Gaussianas} & \textbf{Tempo médio (ms)} & \textbf{p95 (ms)} & \textbf{Hz estimado} & \textbf{Quadros $\leq$ orçamento de 72 Hz} \\ \hline
\textit{baseline} & 195.760 & 28,903 & 31,193 & 34,60 & 0,00\% \\ \hline
\textit{top-k} 100k & 100.000 & 21,541 & 28,717 & 46,42 & 0,00\% \\ \hline
\textit{top-k} 50k & 50.000 & 13,887 & 14,508 & 72,01 & 52,15\% \\ \hline
\textit{splatfacto-big} & 470.962 & 55,449 & 56,922 & 18,03 & 0,00\% \\ \hline
\end{tabularx}
\legend{Fonte: Elaborado pelo autor.}
\end{quadro}

Como fato medido no computador, a poda para 50 mil gaussianas foi a única variante cujo intervalo médio se aproximou do orçamento de 13,889 ms para 72 Hz. Entretanto, somente 52,15\% dos quadros registrados ficaram dentro desse orçamento, e tanto o percentil 95 quanto o máximo o excederam. Na inspeção das capturas estéreo, a cadeira permaneceu reconhecível, enquanto fragmentação e perda de detalhe no piso e no fundo foram mais aparentes com a poda; a variante \textit{splatfacto-big} aparentou maior continuidade visual que as podas, porém obteve o maior tempo de quadro. Essa leitura visual é descritiva e não mede conforto ou latência do streaming.

Esses dados não medem o desempenho do Quest: os tempos e os contadores de processo pertencem ao computador que renderizou a aplicação, e o caminho inclui codificação, transmissão e decodificação de vídeo. Também não houve CSV novo do OVR Metrics nas sessões por Link. Assim, a contingência demonstra o comportamento da representação em hardware de desktop e auxilia o diagnóstico do compromisso entre custo e qualidade, mas não altera a conclusão da avaliação nativa no Meta Quest 3S.

O Quadro \ref{qua:recursos_pcvr} complementa os intervalos da aplicação com os tempos de GPU informados pela Unity e os contadores de memória associados ao processo da aplicação no Windows. As médias de memória foram obtidas das amostras disponíveis durante a coleta externa, cuja janela pode incluir inicialização e captura visual além da janela interna de medição. O conjunto residente corresponde ao contador \textit{Working Set}; o uso dedicado de GPU corresponde ao contador \textit{Dedicated Usage} associado ao processo. Esses contadores não representam a memória do headset nem incluem necessariamente todos os processos envolvidos na transmissão.

\begin{quadro}[H]
\centering
\footnotesize
\caption{Tempos de GPU e contadores de memória do processo de renderização no computador.}
\label{qua:recursos_pcvr}
\renewcommand{\arraystretch}{1.2}
\begin{tabularx}{\textwidth}{|p{3.0cm}|X|X|X|}
\hline
\textbf{Representação} & \textbf{GPU Unity média (ms)} & \textbf{Conjunto residente médio (MiB)} & \textbf{Uso dedicado de GPU médio (MiB)} \\ \hline
\textit{baseline} & 28,601 & 722,04 & 970,54 \\ \hline
\textit{top-k} 100k & 15,752 & 747,47 & 970,76 \\ \hline
\textit{top-k} 50k & 10,568 & 713,20 & 957,11 \\ \hline
\textit{splatfacto-big} & 37,203 & 786,55 & 1.021,43 \\ \hline
\end{tabularx}
\legend{Fonte: Elaborado pelo autor.}
\end{quadro}

A redução do número de gaussianas foi acompanhada por menores tempos médios de GPU nas três variantes principais. Entretanto, os contadores de memória do processo não acompanharam proporcionalmente a redução do PLY. Os dados disponíveis não isolam a contribuição do modelo, dos recursos da Unity e dos demais recursos gráficos; por isso, não permitem atribuir todo o uso de memória às gaussianas.

\subsubsection{Comparação descritiva entre execução nativa e PCVR}

O Quadro \ref{qua:comparacao_nativo_pcvr} e as Figuras \ref{fig:comparacao_tempo_pcvr} e \ref{fig:comparacao_taxa_pcvr} apresentam as mesmas representações nas duas rotas. Para a comparação, foram utilizados o intervalo entre quadros registrado pela aplicação e sua taxa derivada, calculada como mil dividido pelo intervalo médio em milissegundos. Os contadores OVR da execução nativa não foram usados como substitutos desses valores.

\begin{quadro}[H]
\centering
\footnotesize
\caption{Comparação descritiva do intervalo médio entre quadros nas duas rotas.}
\label{qua:comparacao_nativo_pcvr}
\renewcommand{\arraystretch}{1.2}
\begin{tabularx}{\textwidth}{|p{3.0cm}|X|X|X|}
\hline
\textbf{Representação} & \textbf{Quest nativo (ms)} & \textbf{PCVR (ms)} & \textbf{Redução descritiva (\%)} \\ \hline
\textit{baseline} & 45,053 & 28,903 & 35,85 \\ \hline
\textit{top-k} 100k & 26,301 & 21,541 & 18,10 \\ \hline
\textit{top-k} 50k & 17,133 & 13,887 & 18,94 \\ \hline
\textit{splatfacto-big} & 65,625 & 55,449 & 15,51 \\ \hline
\end{tabularx}
\legend{Fonte: Elaborado pelo autor.}
\end{quadro}

\begin{figure}[H]
\centering
\includegraphics[width=\textwidth]{figuras/resultados/quest_pcvr_frame_time_v01.png}
\caption{Intervalo médio entre quadros da aplicação no Quest nativo e no PCVR.}
\label{fig:comparacao_tempo_pcvr}
\legend{Fonte: Elaborado pelo autor.}
\end{figure}

\begin{figure}[H]
\centering
\includegraphics[width=\textwidth]{figuras/resultados/quest_pcvr_frame_rate_v01.png}
\caption{Taxa derivada do intervalo médio entre quadros no Quest nativo e no PCVR.}
\label{fig:comparacao_taxa_pcvr}
\legend{Fonte: Elaborado pelo autor.}
\end{figure}

Na \textit{baseline}, a taxa derivada passou de 22,20 quadros/s na execução nativa para 34,60 quadros/s no PCVR. Para a poda 100k, os valores foram 38,02 e 46,42 quadros/s; para a poda 50k, 58,37 e 72,01 quadros/s. A variante de maior capacidade passou de 15,24 para 18,03 quadros/s. Embora o PCVR tenha apresentado menores intervalos médios em todas as configurações, a ordenação das variantes por custo permaneceu a mesma: a poda 50k foi a mais rápida e a configuração de maior capacidade foi a mais lenta.

O resultado médio de 72,01 quadros/s da poda 50k deve ser interpretado considerando a atualização reportada de 72 Hz nessa sessão. O percentil 95 de 14,508 ms e o máximo de 25,224 ms ultrapassaram o orçamento de 13,889 ms. Portanto, a proximidade da média com a atualização do visor não demonstrou regularidade de todos os quadros. Além disso, os intervalos medidos pela aplicação não quantificam a latência completa entre movimento da cabeça e apresentação da imagem, nem distinguem quadros apresentados pelo compositor de eventuais quadros reprojetados.

A comparação é descritiva: a série Android principal teve cinco execuções por variante, enquanto o PCVR teve uma; a variante de maior capacidade teve um diagnóstico em cada rota. As diferenças também envolvem sistema operacional, API gráfica, hardware e transmissão. Assim, os percentuais apresentados expressam diferenças entre as condições observadas, sem isolar um efeito causal da GPU ou do streaming. A maior capacidade do computador não eliminou a perda de fidelidade das representações podadas nem estabeleceu a adequação da execução nativa no headset.

\subsection{Síntese e limitações dos resultados}

No estudo de caso avaliado, os fatos medidos mostram uma associação entre menos gaussianas, menor tempo de quadro e degradação das métricas offline nas variantes \textit{top-k}. A \textit{baseline} forneceu a referência comparativa; a 50k foi a mais próxima da referência de desempenho entre as variantes da série, mas também teve os menores PSNR e SSIM e o maior LPIPS. A configuração de maior capacidade forneceu indícios qualitativos favoráveis em uma observação, porém teve o maior tempo de quadro no diagnóstico nativo. A interpretação é a existência de um compromisso entre complexidade, desempenho e fidelidade para esse protocolo; ela não demonstra que toda poda produzirá o mesmo efeito em outra cena ou implementação.

As limitações centrais são o uso de uma única cena, a condição formal de pose fixa e a natureza diagnóstica de uma única execução da \textit{splatfacto-big}. A avaliação visual estruturada das podas e da \textit{baseline} é qualitativa; os relatos de travamento e desconforto durante movimento de cabeça são complementares à condição estacionária e podem ter sido influenciados pela execução sequencial das variantes. A contingência PCVR acrescenta somente uma execução por variante, renderizada no computador e transmitida ao headset; ela não reduz essas limitações nem é evidência de viabilidade \textit{standalone}. Esses limites são preservados explicitamente para que a análise não extrapole a evidência disponível.
